% Printed 391, beginning at the dated heading; preceding header is in the prior component.
\begin{center}
\textsc{Gauss} an \textsc{Gerling}.\quad Göttingen, 12.~September 1838.
\end{center}
\par\noindent\hspace*{\parindent}\mbox{.\ .\ .\ .\ .\ .}\quad Was den Längenunterschied zu Paris betrifft, so habe ich selbst die Ergebnisse aus Sternbedeckungen usw., die \textsc{Triesnecker}, \textsc{Wurm} u.~a. gewonnen haben, nie zusammengestellt. Sie selbst sind also ebenso gut wie ich selbst imstande, \textls[120]{dieses} Ergebnis zu ermitteln. Dagegen habe ich aus den Längenunterschieden zwischen Jever und Göttingen und ebenso zwischen Bentheim und Göttingen, wie sie sich einerseits aus meinen trigonometrischen Messungen bzw. denen meines Sohnes ergaben, und aus den Längenunterschieden von Jever und Bentheim gegenüber Paris, wie sie andererseits \textsc{Krayenhoff} angibt, den Längenunterschied zwischen Göttingen und Paris abgeleitet, den Sie in \textsc{Hardings} Ephemeriden angegeben finden. Eine neue Rechnung habe ich deshalb nicht angestellt. Bei einer vollständigen Ausgleichung aller Dreiecke, die das Oldenburgische wie ein Kranz umgeben, ergibt sich für Jever eine, allerdings nur sehr geringe, Änderung. Man könnte Jever auch ganz weglassen und stattdessen Emden und Onstwedde verwenden, bis wohin sich \textsc{Hartmanns} Messungen erstreckt haben. Die Änderung des Endergebnisses wäre jedenfalls ganz unbedeutend, und ich habe es bisher umso weniger der Mühe wert gefunden, deshalb eine neue Rechnung anzustellen, weil \textsc{Krayenhoffs} Messungen, wie Sie meinem Supplem. theor. comb. entnehmen können, doch weit schlechter sind als ihr Ruf.
\par\noindent\hspace*{\parindent}Eine dritte Bestimmung könnten Sie aus Altona gewinnen, dessen Längenunterschied gegenüber Göttingen trigonometrisch bestimmt ist (siehe Breitenunterschied [Art.~19]), ebenso wie der Unterschied gegenüber Greenwich chronometrisch bestimmt ist. Ich selbst habe die betreffenden Daten aber nicht zusammengestellt und muss das daher Ihnen überlassen.\quad\mbox{.\ .\ .\ .\ .\ .}

% Printed 392.
\par\clearpage
\noindent\makebox[0pt][l]{\fontsize{13}{16}\selectfont 392}\hfill{\small\textbf{BRIEFWECHSEL MIT GERLING.}}\hfill\null\par\medskip
\begin{center}
\textsc{Gauss} an \textsc{Gerling}.\quad Göttingen, 14.~November 1838.
\end{center}
\par\noindent\hspace*{\parindent}\mbox{.\ .\ .\ .\ .\ .}\quad Dass ich auf Ihren vorletzten Brief noch nicht weiter geantwortet habe, müssen Sie gütigst damit entschuldigen, dass ich dem in meinem letzten Brief [vom 12.~September] Mitgeteilten eigentlich kaum noch etwas hinzuzufügen wusste.
\par\noindent\hspace*{\parindent}An Tatsachen kann ich wirklich nur noch einen Hinweis geben. In \textsc{Lindenau-Bohnenbergers} Zeitschrift, Band IV, S.~119, finden Sie die relative Lage der alten und der neuen hiesigen Sternwarte bereits angegeben. Allenfalls kann ich die folgenden beiden Anmerkungen hinzufügen.
\par\noindent\hspace*{\parindent}1) Der Längenunterschied ist in \textls[120]{Zeit} angegeben.
\par\noindent\hspace*{\parindent}2) Die Position der neuen Sternwarte an der angegebenen Stelle bezieht sich auf das Zentrum der Rotunde, dessen Lage gegenüber der Mitte der Achse des \textsc{Reichenbach}schen Meridiankreises ich Ihnen, wenn ich mich nicht irre, bereits früher mitgeteilt habe. Wenn Sie aber die relative Lage des letzteren Punktes gegenüber dem Zentrum der alten Sternwarte selbst genau berechnen wollen, schreibe ich Ihnen dazu aus meinem Koordinatenverzeichnis Folgendes ab:
\begin{center}
\setlength{\tabcolsep}{4pt}
\begin{tabular}{@{}lrr@{}}
Zentrum der Kuppel & \(-3{,}104\) & \(-7{,}324\)\\
Mitte der Achse des \textsc{Reichenbach}schen Meridiankreises & \(0\) & \(0\)\\
Mitte der alten Sternwarte & \(-193{,}54\) & \(+541{,}9\).
\end{tabular}
\end{center}
\par\noindent\hspace*{\parindent}Die Einheit ist das Meter. Bei der ersten Zahl bedeutet \(+\) südlich, bei der zweiten \(+\) westlich.
\par\noindent\hspace*{\parindent}Wenn Sie nun \textls[120]{meine} Meinung über die \textls[120]{Reduktionszahl} wissen wollen, die Sie bei der Angabe Ihrer Längen gegenüber Paris verwenden sollen, so kann ich nur sagen, dass ich das für ganz unerheblich halte, sofern Sie mir zugleich mitteilen, \textls[120]{welche} Zahl Sie zugrunde gelegt haben. In der Tat bleibt die Wahlfreiheit in dieser Hinsicht, oder die Größe des Spielraums, unabhängig davon, was Sie wählen,
\par\noindent\hspace*{\parindent}\hangindent=2.8em\hangafter=1 A.\quad \textls[120]{höchst} unbedeutend für jeden praktischen Zweck, z.~B. für das Zeichnen von Karten.
\par\noindent\hspace*{\parindent}\hangindent=2.8em\hangafter=1 B.\quad zwanzig- oder fünfzigmal größer als alle \textls[120]{relativen} Differenzen zwischen all Ihren Ansätzen.

% Printed 393.
\par\clearpage
\noindent\null\hfill{\small\textbf{ZUR HANNOVERSCHEN TRIANGULATION.}}\hfill\makebox[0pt][r]{\fontsize{13}{16}\selectfont 393}\par\medskip
\par\noindent\hspace*{\parindent}Da nun die letzteren eigentlich das \textls[120]{Wesen Ihrer Arbeit} ausmachen, scheint mir damit mein oben ausgesprochenes Urteil hinreichend gerechtfertigt.
\par\noindent\hspace*{\parindent}In Ihrem vorletzten Brief bezeichnen Sie die Abweichungen des auf verschiedenen Wegen bestimmten Längenunterschieds zwischen Paris und Göttingen, die bis zu \(\frac{3}{4}\) Zeitsekunden betragen, als sehr beträchtlich. Ich würde sie nicht so nennen.
\par\noindent\hspace*{\parindent}In der Tat scheinen Sie mir dabei ganz außer Acht gelassen zu haben, was ich in meiner Schrift über den Breitenunterschied zwischen Göttingen und Altona auf S.~73 [*)] entwickelt habe. Wenn man auf geodätischem Weg zwischen zwei nicht allzu weit auseinanderliegenden Punkten über zwei verschiedene Routen Längenunterschiede fände, die um 11 Bogensekunden voneinander abwichen, wäre das ein ganz ungeheurer Fehler. Er würde beweisen, dass ganz enorm schlechte Beobachtungen eingeflossen wären oder dass schlecht gerechnet worden wäre.
\par\noindent\hspace*{\parindent}Aber Paris und Göttingen liegen doch sehr weit auseinander. \textsc{Delambres} und \textsc{Krayenhoffs} Messungen sind, glaube ich, an Genauigkeit mit den unseren gar nicht zu vergleichen und auch nach ganz anderen Methoden und mit ganz anderen Elementen berechnet. Ich würde mich daher nicht wundern, wenn der so gefundene Längenunterschied um einige Bogensekunden von einem anderen abwiche, der ebenfalls geodätisch, aber z.~B. auf einem Weg über Strassburg, Mannheim usw. bestimmt wäre.
\par\noindent\hspace*{\parindent}Allein \textls[120]{davon} ist hier gar nicht die Rede, sondern vom Unterschied der astronomisch bestimmten Längen, die, genau betrachtet, etwas \textls[120]{ganz anderes} bedeuten als die geodätischen. Jene beziehen sich auf die fortschreitende Änderung der Lage der Lotlinien und der durch sie parallel zur Erdachse gelegten Ebenen, die anderen auf die Entfernungen auf der Erdoberfläche, die \textls[120]{gar kein Ellipsoid} ist, sondern für die ein Ellipsoid nur als eine Art Näherung angesehen werden kann.
\par\noindent\hspace*{\parindent}Angesichts der unzähligen bekannten Abweichungen zwischen geodätischer und astronomischer Bestimmung der Breite muss man bei der Länge überall mit Abweichungen derselben Größenordnung rechnen. Eine Differenz
\par\vfill
\noindent\rule{.22\linewidth}{.4pt}\par\smallskip
{\footnotesize\noindent\hspace*{\parindent}[*) S.~49/50 dieses Bandes.]\par}
\medskip\noindent{\small IX.}\hfill 50\hspace*{1.5em}\par

% Printed 394.
\clearpage
\noindent\makebox[0pt][l]{\fontsize{13}{16}\selectfont 394}\hfill{\footnotesize\textbf{BRIEFWECHSEL MIT GERLING, ZUR HANNOVERSCHEN TRIANGULATION.}}\hfill\null\par\medskip
\noindent von \(11^{\prime\prime}\), selbst zwischen zwei Orten, die einander \textls[120]{viel} näher lägen als Göttingen und Paris, ist im Grunde gar nichts. Es wäre überhaupt nichts Besonderes, wenn sich zwischen Göttingen und Marburg sogar eine doppelt so große fände.
\par\noindent\hspace*{\parindent}Übersehen Sie also nicht, dass die von Ihnen gefundenen Unterschiede selbst dann sehr wohl vollkommen genau richtig sein könnten, wenn \textls[120]{alle} Beobachtungen vollkommen und absolut fehlerfrei wären (was sie nicht sind). Es gibt \textls[120]{gar keinen} Grund zu erwarten, dass
\par\noindent\hspace*{\parindent}1) der durch geodätische Messungen unter der Annahme bestimmter Ellipsoiddimensionen berechnete Unterschied zwischen Göttingen und Paris,
\par\noindent\hspace*{\parindent}2) der durch astronomische Beobachtungen gefundene,
\par\noindent\hspace*{\parindent}3) der durch eine Kombination gefundene, nämlich geodätisch zwischen Göttingen und Altona und astronomisch zwischen Altona-Greenwich-Paris,
\par\noindent miteinander übereinstimmen, da diese drei Zahlen tatsächlich ganz verschiedene Dinge bedeuten.
\par\noindent\hspace*{\parindent}Wenn ich bei einer künftigen Veröffentlichung der Ergebnisse meiner Messungen einen Teil davon in Form von Breiten- und Längenunterschieden angeben sollte, werde ich die letzteren ganz gewiss nicht auf Paris beziehen, das mich gar nichts angeht. Wenn Sie es aber aus irgendeinem Grund für ratsam halten, die Ihren in einer solchen Form mitzuteilen, wiederhole ich, dass ich es für ganz gleichgültig halte, welche Reduktionszahl Sie verwenden, sofern für den Leser klar bleibt, dass die darin enthaltenen Sekundenbruchteile für sich genommen gar keinen Anspruch auf eine genaue Bedeutung erheben können. Sie dienen nur dazu, eine beliebige Differenz zwischen zwei Ihrer Zahlen mit einer solchen Genauigkeit wiederzugewinnen, dass nichts von der Genauigkeit Ihrer Messungen selbst verloren geht. Ohne Ihnen vorgreifen zu wollen, würde ich unter der eben genannten Voraussetzung die Zahl zugrunde legen, die sich taliter qualiter aus der vorhandenen geodätischen Verbindung zwischen Göttingen und Paris ergeben hat, da die astronomischen Längen, die etwas ganz anderes bedeuten, hier eigentlich gar nicht infrage kommen.\quad\mbox{.\ .\ .\ .\ .\ .}
\begin{center}\rule{.12\linewidth}{.4pt}\end{center}

% Printed 395.
\par\clearpage
\noindent\null\hfill{\footnotesize\textbf{BEMERKUNGEN. BRIEFWECHSEL ZUR HANNOVERSCHEN TRIANGULATION.}}\hfill\makebox[0pt][r]{\fontsize{13}{16}\selectfont 395}\par\medskip
\begin{center}\textbf{BEMERKUNGEN.}\end{center}
\begingroup\small
\par\noindent\hspace*{\parindent}Die vorstehenden Briefe wurden nach den Originalen abgedruckt.
\par\noindent\hspace*{\parindent}Dem Brief an \textsc{Bessel} vom 26.~Dezember 1821, S.~349/351, war eine Dreiecksskizze beigefügt, die die 5 Dreiecke von der Seite Hohehagen-Inselsberg bis zur Seite Hils-Lichtenberg (vgl. die Figur auf S.~299) enthält und außerdem die Visuren vom Hils zum Deister und zum Brelingerberg angibt. Der Brelingerberg, nördlich von Hannover, und der Wohlenberg, nordwestlich von Braunschweig, waren 1821 bei der Erkundung als Dreieckspunkte vorgesehen. Der erste erwies sich jedoch zu Beginn der Arbeiten im folgenden Jahr als ungeeignet, der zweite wurde als überflüssig weggelassen. Wahrscheinlich gehörte auch zu dem Brief an \textsc{Bessel} vom 15.~November 1822, S.~351/358, eine Karte.
\par\noindent\hspace*{\parindent}Zu den Briefen an \textsc{Olbers} vom 4. und 8.~Juli 1824, S.~369/372, ist noch anzumerken, dass der darin erwähnte Punkt Wentel nicht in das System aufgenommen wurde.
\par\noindent\hspace*{\parindent}Es war, wie es im Arbeitsbericht von \textsc{Gauss} für 1824 an das hannoversche Ministerium heißt, »schneller als man bei den großen Schwierigkeiten hatte hoffen können, ein sehr gutes Dreieckssystem bis Bremen gebildet.« »Doch umso trüber wurden nun die Aussichten für die weitere Fortsetzung der Dreiecke.\ \mbox{.\ .\ .\ .}\quad Nach dem Fehlschlagen aller Versuche, einen neuen Dreieckspunkt an Brüttendorf anzuschließen, blieb noch ein Behelf übrig, nämlich den Steinberg in das Dreieckssystem aufzunehmen.« Aber auch an die Seite Bremen-Steinberg ließ sich kein neues Dreieck anschließen. »Ohne einen gerade jetzt eingetretenen glücklichen Umstand hätte es um die Fortsetzung der Messungen sehr schlecht gestanden, da nun alle Möglichkeiten ausgeschöpft schienen.« Der erwähnte glückliche Umstand war die Entdeckung von \textsc{Gauss} auf dem Ansgariusthurm in Bremen, dass die Spitze des Turms von Zeven von dort aus gerade noch zu sehen war. Dass Bremen von Zeven aus sichtbar war, war früher wegen des Moorrauchs nicht bemerkt worden. Zeven ließ sich außer mit Bremen auch mit den vorhergehenden Punkten Steinberg, Wilsede und Litberg verbinden. Von Bremen-Zeven aus konnte die Dreieckskette dann fortgesetzt werden. (Vgl. dazu den später folgenden Auszug aus dem Arbeitsbericht von \textsc{Gauss} für 1825 sowie den ebenfalls noch folgenden Brief an \textsc{Olbers} aus dem Juli 1825.)
\par\noindent\hspace*{\parindent}Die im Brief an \textsc{Gerling} vom 11.~August 1823, S.~385/386, erwähnten, mit dem Heliotrop übermittelten Zahlenzeichen hatten (nach einigen kurzen Notizen in Beobachtungs- und Rechnungsheften zur Gradmessung) vorher vereinbarte, teilweise von Station zu Station wechselnde Bedeutungen (vgl. auch den Anfang des Briefs an \textsc{Gerling} vom 1.~September 1823, S.~387). Nach einer Aufzeichnung bedeutete z.~B. 1: Es geht so gut, 2: Das Licht soll abgeschwächt werden, 3: Für heute soll aufgehört werden, 1\,.\,1: Aufbruch zum nächsten Punkt.
\par\noindent\hspace*{\parindent}Auf die Gradmessungsarbeiten in den Jahren 1821/1825 beziehen sich auch die Briefe an \textsc{Schumacher} vom 6. und 30.~Mai, 11.~Juli, 29.~September und 24.~Oktober 1821; vom 10.~Mai, 10.~Juni, 6. und 30.~August, 6., 18., 24. und 29.~September, 8.~Oktober und 10.~November 1822; vom 8. und 18.~Juni, 23.~Juli, Anfang August, 21.~August, 18.~September und 23.~Oktober 1823; vom 24.~Juni, 1.~Juli, 27.~September, 17.~Oktober, 28.~November 1824 und vom Januar 1825; vom 29.~April, 14., 20. und 26.~Juni, 11. und 15.~Juli und 14.~August 1825. (Vgl. den Briefwechsel \textsc{Gauss-Schumacher}, Erster Band, S.~229/245, S.~265/291, S.~311/325, S.~334/336, S.~397/413, S.~426/430. Zweiter Band, S.~1/3, S.~14/31.)
\par\noindent\hspace*{\parindent}Die mehrfach erwähnte Triangulation des Kurfürstentums Hannover durch französische Ingenieurgeografen unter der Leitung von Oberstleutnant \textsc{Epailly} fand 1804 und 1805 während der französischen Besatzung statt. Sie sollte die Unterlagen für die kartografische Aufnahme liefern. Die Tabelle der \textsc{Epailly}schen Dreiecke nebst einem Bericht \textsc{Epaillys} über seine Arbeiten erhielt \textsc{Gauss} Anfang 1821 durch \textsc{Laplaces} Vermittlung vom französischen Kriegsministerium.
\par\noindent\hfill\textsc{Krüger}.\par
\endgroup
\begin{center}\rule{.12\linewidth}{.4pt}\end{center}
\vfill\noindent\hfill 50\textsuperscript{*}\hspace*{1.5em}\par

% Printed 396.
\clearpage
\noindent\makebox[0pt][l]{\fontsize{13}{16}\selectfont 396}\hfill{\small\textbf{ZUR HANNOVERSCHEN TRIANGULATION.}}\hfill\null\par\medskip
\begin{center}
{\large\textbf{VERÖFFENTLICHUNGEN}}\par\bigskip
{\large\textbf{ZUR HANNOVERSCHEN TRIANGULATION.}}\par\medskip
[II.]\par\bigskip
\rule{.74\linewidth}{.4pt}\\[-.5ex]
{\small Astronomische Nachrichten, Band I, Nr.~7, Februar 1822, S.~105---106.}\\[-.5ex]
\rule{.74\linewidth}{.4pt}
\end{center}
\par\noindent\hspace*{\parindent}\{Es ist bekannt, dass die hannoversche Regierung die von Seiner Majestät dem König von Dänemark, diesem erhabenen Förderer der Wissenschaften, begonnene Gradmessung, die sich in der Breite vom nördlichsten Punkt Jütlands bis zur südlichsten Grenze von Lauenburg erstreckt, fortsetzt und die Ausführung dieser Arbeit Herrn Hofrat \textsc{Gauss} in Göttingen übertragen hat. Unsere gemeinsamen Dreiecke sind Hamburg-Hohenhorn-Lüneburg und Hohenhorn-Lauenburg-Lüneburg. In diesen Dreiecken habe ich den Winkel in Hamburg allein beobachtet, die Winkel in Lüneburg haben Herr Hofrat \textsc{Gauss} und ich gemeinsam beobachtet (der Unterschied zwischen unseren Messungen lag unter \(0^{\prime\prime}{,}3\)). Die Winkel in Hohenhorn und Lauenburg wurden zur selben Zeit (1818) von Herrn Kapitän v.~\textsc{Caroc} gemessen. Im vergangenen Jahr hat der Herr Hofrat seine Dreiecke an der südlichen Grenze begonnen. Der Auszug aus seinem Brief, den ich hier mitteile, enthält einen kurzen Bericht darüber.
\par\noindent\hfill\textit{S}[\textit{chumacher}].\}\par
\begin{center}Auszug aus einem Brief des Herrn Hofrats und Ritters \textsc{Gauss}.\end{center}
\par\noindent\hspace*{\parindent}Ich habe das Vergnügen, Ihnen den Stand meiner Triangulation auf der beiliegenden kleinen Karte mitzuteilen. Die kräftigen Linien zeigen Richtungen an, die auf bereits vorgenommenen Messungen beruhen, die punktierten zeigen geplante Richtungen an. Ich habe leider Grund zu befürchten, dass der Brelingerberg weder vom Wohlenberg noch vom Lichtenberg aus sichtbar ist (ich selbst war an keinem der 3 Orte). Überhaupt wird es sehr schwierig sein, in der Lüneburger Heide große Dreiecke zu bilden.

% Printed 397.
\par\clearpage
\noindent\null\hfill{\small\textbf{ZUR HANNOVERSCHEN TRIANGULATION.}}\hfill\makebox[0pt][r]{\fontsize{13}{16}\selectfont 397}\par\medskip
\par\noindent\hspace*{\parindent}Brelingerberg, Deister, Lichtenberg und Inselsberg wurden durch Heliotroplicht sichtbar gemacht, auf der Brockenstation auch Hohehagen und Hils, da die dort errichteten Signale nur selten (das letztere nur wenige Minuten), ich will nicht sagen zu beobachten, sondern auch nur zu \textls[120]{sehen} waren. Mit dem Heliotrop fällt jede Schwierigkeit weg.
\par\noindent\hspace*{\parindent}Ich lasse jetzt noch zwei weitere anfertigen (nach der neuen Bauart), von denen eines bald fertig sein wird. Dass meine mit dem Sextanten-Heliotrop auf dem Brocken gegebenen Telegrafenzeichen auf dem Hohehagen (Entfernung \(70000\) Meter \(=9\frac{1}{3}\) geografische Meilen) \textls[120]{mit bloßen Augen} gesehen wurden, habe ich Ihnen, glaube ich, bereits in meinem letzten Brief mitgeteilt. Das bisherige Heliotrop kann mein ältester Sohn Joseph schon recht gut einstellen und ausrichten. Mit dem neuen wird es wohl noch etwas leichter gehen.
\par\noindent\hspace*{\parindent}Die Richtung vom Hils auf Hannover ist zwar ebenfalls mit größter Genauigkeit gemessen worden. Hannover wird aber vermutlich kein Hauptdreieckspunkt werden, da man von dort nach NO nur eine sehr begrenzte Aussicht hat und insbesondere den Wohlenberg nicht sehen kann. Auch wären an dem dortigen Turm erst große Änderungen nötig, wenn dort ein Theodolit aufgestellt werden sollte. Der Deister wird nach allen Richtungen noch eine weite Aussicht bieten, und vielleicht wird von dort sogar noch der Falkenberg zu sehen sein.
\begin{center}
\rule{.74\linewidth}{.4pt}\\[-.5ex]
{\small Astronomische Nachrichten, Band I, Nr.~24, Dezember 1822, S.~441---444.}\\[-.5ex]
\rule{.74\linewidth}{.4pt}\par\medskip
Auszug aus einem Schreiben des Herrn Hofrats \textsc{Gauss} an den Herausgeber.\\
Göttingen 1822, Nov.~10.
\end{center}
\par\noindent\hspace*{\parindent}Da Sie im \(7.\) Heft der Astronomischen Nachrichten eine Mitteilung über den Stand meiner Triangulation am Ende des Jahres 1821 gebracht haben, versäume ich es nicht, Ihnen einen kurzen Bericht über den gegenwärtigen Stand der Arbeiten zu schicken.
\par\noindent\hspace*{\parindent}Im vergangenen Jahr wurden die fünf Stationen Göttingen, Meridianzeichen, Hohehagen, Hils und Brocken abgeschlossen und vier Punkte für die weitere Fort-

% Printed 398.
\par\clearpage
\noindent\makebox[0pt][l]{\fontsize{13}{16}\selectfont 398}\hfill{\small\textbf{ZUR HANNOVERSCHEN TRIANGULATION.}}\hfill\null\par\medskip
\noindent setzung der Arbeiten markiert, nämlich Lichtenberg, Deister, Wohlenberg und Brelingerberg. Ich begann die Arbeiten des laufenden Jahres mit einer Erkundungsreise in die Lüneburger Heide. Ich hielt sie umso mehr für notwendig, als ich die großen Schwierigkeiten, in diesem flachen Land, das keine nennenswerten Anhöhen hat und überall schachbrettartig mit Wald bedeckt ist, ein Dreiecksnetz zu bilden, bereits aus den Berichten des Obersten \textsc{Epailly} kannte. Dieser hatte die Schwierigkeiten in den Jahren 1804 und 1805 unüberwindlich gefunden und deshalb die Verbindung zwischen Hamburg und dem südlichen Teil Hannovers mittels einer Reihe von Dreiecken entlang der Weser bis zu ihrer Mündung und dann wieder die Elbe hinauf hergestellt.
\par\noindent\hspace*{\parindent}Ich fand den Brelingerberg, der 1821 vom Hils aus angepeilt worden war, \textls[120]{unbrauchbar}, da er sich nicht mit Lichtenberg und dem Wohlenberg verbinden ließ, aber ebenso wie den Wohlenberg auch \textls[120]{überflüssig}, da sich sowohl der Platz bei Garssen als auch der Falkenberg unmittelbar mit dem Lichtenberg verbinden ließen. Ich übergehe die großen Schwierigkeiten, mit denen ich zu kämpfen hatte, um die Dreiecke von Garssen und Falkenberg aus weiterzuführen. Diese Schwierigkeiten sind jetzt überwunden, und das Netz bietet durch seinen Aufbau zahlreiche Kontrollen, die zu meiner größten Zufriedenheit ausgefallen sind. Ich bemerke nur, dass alle meine Dreieckspunkte auf ebener Erde liegen. Ein etwa \(3\frac{1}{2}\)---4 Fuß hohes gemauertes Steinpostament dient zur Aufstellung des Heliotrops und des Theodoliten. Mehrere Linien, namentlich die von Falkenberg nach Wilsede, von Hauselberg nach Breithorn, von Breithorn nach Scharnhorst und von Scharnhorst [nach Garssen], erforderten beträchtliche Schneisen durch den Wald. Die genaue Vorausbestimmung der Richtung dieser Schneisen erforderte aufwendige Vorbereitungen.
\par\noindent\hspace*{\parindent}Im Laufe des Sommers habe ich die Stationen Lichtenberg, Deister, Garssen, Falkenberg, Hauselberg, Breithorn, Wulfsode, Wilsede und Scharnhorst vollständig abgeschlossen und auf dem Timpenberg auch die betreffenden Winkel vorläufig gemessen. Damit ist Hamburg schon vorläufig angeschlossen, ebenso Lüneburg, da Sie den Winkel zwischen Wilsede und Lüneburg auf dem Michaelisthurm in Hamburg vorläufig gemessen haben. Hier einige vorläufige Ergebnisse, deren absoluter Maßstab vorläufig auf der von \textsc{Zach}schen Basis bei

% Printed 399.
\par\clearpage
\noindent\null\hfill{\small\textbf{ZUR HANNOVERSCHEN TRIANGULATION.}}\hfill\makebox[0pt][r]{\fontsize{13}{16}\selectfont 399}\par\medskip
\noindent Gotha beruht, an die ich mich mittels der Seite vom Inselsberg zum Brocken angeschlossen habe.
\begin{center}
\setlength{\tabcolsep}{3pt}
\begin{tabular}{@{}lrrr@{\qquad}rrrl@{}}
 & \multicolumn{3}{c}{\kern-10pt Breite\kern10pt} & \multicolumn{4}{c}{\kern-5pt\shortstack{Länge von\\Göttingen}\kern5pt}\\[-1.5ex]
 & \multicolumn{3}{c}{\kern-10pt\(\overbrace{\hphantom{53\quad33\quad31{,}4}}\)\kern10pt} & \multicolumn{4}{c}{\kern-5pt\(\overbrace{\hphantom{0\quad27\quad29{,}5\quad\text{östl.}}}\)\kern5pt}\\[-1.2ex]
Hamburg, Michaelisthurm & \(53^\circ\) & \(33'\) & \(1^{\prime\prime}{,}8\) & \(0^\circ\) & \(2'\) & \(3^{\prime\prime}{,}0\) & östl.\\
Lüneburg, Michaelisthurm & 53 & 15 & \(5{,}5\) & 0 & 27 & \(29{,}5\) & \multicolumn{1}{c}{»}\\
Celle, südl. Schlossturm & 52 & 37 & \(31{,}4\) & 0 & 8 & \(4{,}9\) & \multicolumn{1}{c}{»}\\
Göttinger Sternwarte, Platz & & & & & & &\\
\quad des \textsc{Reichenbach}schen & & & & & & &\\
\quad Meridiankreises & 51 & 31 & \(48{,}7\) & & 0 & &
\end{tabular}
\end{center}
\par\noindent\hspace*{\parindent}Die Orientierung meines Dreieckssystems geht von meinem Meridianzeichen aus. Auf Hamburg übertragen, weicht sie von den Azimuten, die Sie mir mitgeteilt haben, nur um \(1^{\prime\prime}{,}4\) ab. Um die absoluten Längen der Linien genauer zu bestimmen, warte ich nur noch auf die Mitteilung der Länge Ihrer Basis und der Dreiecke, die sie mit Ihren Hauptpunkten verbinden. Die eine meiner Dreiecksseiten, Breithorn-Scharnhorst, ließe sich, wie es scheint, ohne unüberwindliche Schwierigkeiten unmittelbar messen.
\par\noindent\hspace*{\parindent}Um meine Dreiecke möglichst zweckmäßig mit den Ihren zu verbinden, hatte ich gewünscht und gehofft, Timpenberg unmittelbar mit Lüneburg verbinden zu können. Man versuchte, eine Schneise zu schlagen. Nachdem sie ein beträchtliches Stück vorangetrieben worden war, zeigte sich aber, dass das dazwischenliegende Gelände nicht tief genug lag, sodass diese unmittelbare Verbindung aufgegeben werden musste. Von Wilsede aus wurde jedoch noch ein Punkt festgelegt, der sich unmittelbar mit Hamburg, Lüneburg und Lauenburg und höchstwahrscheinlich mittels einer Schneise mit Timpenberg verbinden lässt. Alles Weitere muss den Arbeiten des kommenden Jahres vorbehalten bleiben. Von den großen Schwierigkeiten, in einem solchen bewaldeten, flachen Gelände zu arbeiten, hat niemand eine Vorstellung, der nicht unter ähnlichen Umständen gearbeitet hat.
\par\noindent\hspace*{\parindent}Die beiliegende Karte, die im Maßstab \(\frac{1}{2000000}\) gezeichnet ist, wird Ihnen eine anschaulichere Vorstellung vom Erreichten geben. Erst nachdem die übrigen Arbeiten abgeschlossen waren, zeigte sich, dass der Punkt Scharnhorst mithilfe zweier nicht sehr schwierig anzulegender Schneisen unmittelbar

% Printed 400.
\par\clearpage
\noindent\makebox[0pt][l]{\fontsize{13}{16}\selectfont 400}\hfill{\small\textbf{ZUR HANNOVERSCHEN TRIANGULATION.}}\hfill\null\par\medskip
\noindent mit Lichtenberg und dem Deister verbunden werden könnte. Hätte man diesen Platz früher ausfindig machen und seine Eignung und Lage feststellen können, hätte Garssen ganz entfallen können. Vielleicht hole ich im kommenden Jahr die Messung der Winkel des Dreiecks Scharnhorst-Deister-Lichtenberg noch nach.
\par\noindent\hspace*{\parindent}Ich hatte in diesem Jahr außer dem 1821 verwendeten Heliotrop noch zwei weitere neuer Bauart im Einsatz und daneben noch ein anderes Heliotropgerät, das ich immer bei mir trug, um meinen Gehilfen telegrafische Anweisungen zu geben. Ihnen gegenüber ist die Bemerkung überflüssig, dass der von Herrn~\textsc{Schubach} im Astron. Jahrbuch für 1825 gegebene Bericht über die Bauart der Heliotrope ganz auf einem Irrtum beruht und mit \textls[120]{meinen} Heliotropen gar nichts gemein hat. Herr \textsc{Rumpf} hat bereits sieben Heliotrope angefertigt, von denen zwei für die preußische und zwei für die hessische Triangulation bestimmt sind. Zeichnungen beider Bauarten stehen Ihnen auf Wunsch zur Verfügung.
\begin{center}\rule{.12\linewidth}{.4pt}\end{center}
\begin{center}\textbf{BEMERKUNGEN.}\end{center}
\begingroup\small
\par\noindent\hspace*{\parindent}Die dem ersten Artikel in Band I, Nr.~7, der Astronomischen Nachrichten beigegebene Dreieckskarte zeigt in kräftigen Linien im Süden die Dreiecke von der Seite Hohehagen-Inselsberg bis zur Seite Hils-Lichtenberg (vgl. die Figur auf S.~299) sowie die Visuren vom Hils zum Deister und zum Brelingerberg und im Norden die beiden Dreiecke Lüneburg-Hamburg-Hohenhorn und Lüneburg-Hohenhorn-Lauenburg. Punktierte Linien zeigen das Dreieck Lichtenberg-Deister-Wohlenberg, die Richtungen zwischen den Punkten Deister, Brelingerberg, Falkenberg und Wilsede sowie die Richtungen Wilsede-Lüneburg und Wilsede-Hamburg. Wohl aus Versehen fehlt auf dieser Karte die Linie Hils-Nördl. Meridianzeichen.
\par\noindent\hspace*{\parindent}Die dem zweiten Artikel beigefügte Dreieckskarte zeigt die Dreiecke der eigentlichen hannoverschen Gradmessung, also die in der Figur auf S.~299 dargestellten Dreiecke vom Inselsberg bis Hamburg bis zu den Seiten Falkenberg-Wilsede und Wilsede-Hamburg. Auf ihr fehlen allerdings noch der Punkt Nindorf und seine Verbindungen. Dagegen enthält sie die Anschlussdreiecke zwischen den Punkten Lüneburg, Wilsede, Hamburg und Hohenhorn und die beiden dänischen Punkte Syk und Lauenburg.
\par\noindent\hspace*{\parindent}Außer diesen beiden Mitteilungen über die hannoversche Triangulation hat \textsc{Gauss} dazu noch 2 Aufsätze aus dem Jahr 1821 über den »Heliotrop« (die später folgen) und einen Artikel aus dem Jahr 1823 veröffentlicht, »Beobachtete und berechnete Triangulirung im Hannoverschen, Braunschweigschen und Lüneburgschen« (der ebenfalls noch folgt).
\par\noindent\hfill\textsc{Krüger}.\par
\endgroup
\begin{center}\rule{.16\linewidth}{.4pt}\end{center}
% End of printed 400. Printed 401 opens III. NACHLASS ZUR HANNOVERSCHEN TRIANGULATION.
